% Shared preamble for the attribution infographics (iN-*.tex).
% tools/tikz-build.sh compiles each source with \variant set to wide or narrow.
\usepackage[T1]{fontenc}
\usepackage[default,semibold]{sourcesanspro}
\usepackage[italic]{mathastext}
% Larger scripts than the TeX default, so subscripts stay legible at phone width.
\DeclareMathSizes{9}{9}{8}{7}
\DeclareMathSizes{10}{10}{9}{8}
\DeclareMathSizes{11}{11}{9}{8}
\DeclareMathSizes{11.5}{11.5}{9.5}{8}
\usetikzlibrary{arrows.meta,positioning,calc,fit,backgrounds,decorations.pathreplacing,patterns}

\ifdefined\variant\else\def\variant{wide}\fi
\newif\ifnarrow
\def\narrowname{narrow}
\ifx\variant\narrowname\narrowtrue\fi
% The narrow variant is about 320pt wide and shown at about 350px, the wide one
% about 540pt wide and shown at up to 720px: the narrow one uses larger type.
\ifnarrow
  \def\basefont{\fontsize{11}{13}\selectfont}\def\smallfont{\fontsize{10}{12}\selectfont}
  \def\titlefont{\fontsize{11.5}{13.5}\selectfont\bfseries}
\else
  \def\basefont{\fontsize{10}{12}\selectfont}\def\smallfont{\fontsize{9}{11}\selectfont}
  \def\titlefont{\fontsize{11}{13}\selectfont\bfseries}
\fi

% Light-theme page tokens (web/src/styles/_variables.css, embeds.css), fixed:
% every infographic sits on its own light card, so it reads on both themes.
\definecolor{ink}{HTML}{262626}      % --text-color on white
\definecolor{muted}{HTML}{666666}    % --muted-color on white
\definecolor{rule}{HTML}{DCDCDC}     % --border-color, slightly stronger
\definecolor{surface}{HTML}{F5F5F5}  % --neutral-200
\definecolor{heading}{HTML}{904B20}  % h2 colour
\definecolor{info}{HTML}{187C9B}     % embeds --info-color
\definecolor{primary}{HTML}{C65928}  % embeds --primary-color
\definecolor{success}{HTML}{35814C}  % embeds --success-color
\definecolor{danger}{HTML}{BD3F50}   % f13 spread cost
\definecolor{residual}{HTML}{A77913} % f13 residual

\tikzset{
  every picture/.style={color=ink, line width=0.6pt, >={Stealth[length=5pt,width=4pt]},
    font=\basefont},
  small/.style={font=\smallfont},
  title/.style={font=\titlefont},
  box/.style={draw=rule, fill=white, rounded corners=4pt, align=center, inner sep=5pt},
  tag/.style={small, text=muted},
}

% No hyphenation in the narrow text columns of the diagrams.
\hyphenpenalty=10000 \exhyphenpenalty=10000

% Light card behind the whole drawing, below the background layer.
\pgfdeclarelayer{card}
\pgfsetlayers{card,background,main}
\newcommand\drawcard{%
  \begin{pgfonlayer}{card}
    \fill[white, draw=rule, rounded corners=8pt]
      ($(current bounding box.south west)+(-10pt,-10pt)$)
      rectangle ($(current bounding box.north east)+(10pt,10pt)$);
  \end{pgfonlayer}}
